% Original Prism Database Issue skeleton for Nucleic Acids Research.
% Uses the unmodified OUP 1.5 class and its Modern Large page design.
% The sections below are writing suggestions, not an official fixed outline.
\documentclass[unnumsec,webpdf,modern,large]{oup-authoring-template}
\usepackage{booktabs,tabularx,siunitx}
\newcommand{\societylogo}{}
\sisetup{uncertainty-mode=separate}

\begin{document}
\journaltitle{Nucleic Acids Research}
\DOI{}
\copyrightyear{}
\pubyear{}
\access{Draft manuscript}
\appnotes{Database Issue}
\firstpage{1}

\title[Database name and short title]{Database name: concise description of the resource}
\author[1,$\ast$]{First Author}
\author[1,2]{Second Author}
\address[1]{\orgdiv{Department}, \orgname{First institution},
  \orgaddress{\state{City}, \postcode{Postal code}, \country{Country}}}
\address[2]{\orgdiv{Department}, \orgname{Second institution},
  \orgaddress{\state{City}, \postcode{Postal code}, \country{Country}}}
\corresp[$\ast$]{Corresponding author.
  \href{mailto:first.author@example.org}{first.author@example.org}}

% One English paragraph, at most 200 words; include the actual database URL.
% example.org is reserved demonstration content, not a functioning database.
\abstract{Replace this paragraph with a factual description of the
database, its biological scope, data coverage, curation and access
functions. Explain the verified benefit to users or the substantive
changes in an updated release. State the actual resource URL in the
abstract. The URL \url{https://example.org/database} is a reserved example
that must be replaced; it does not identify a database. This starter
reports no collected data, functionality or scientific result.}
\maketitle
\thispagestyle{plain}
\pagestyle{plain}

\section{Introduction}\label{sec:introduction}
Explain the biological need, intended users and relationship to existing
resources. For an update, identify the previous publication and describe
the substantive new content or capabilities. The public reference below
tests citation mechanics and should be retained only if relevant to the
resource~\cite{wilkinson2016fair}.

\section{Materials and Methods}\label{sec:methods}
\subsection{Data sources and provenance}
Describe actual data sources, accessions, release dates, inclusion rules
and permissions. Distinguish imported records, manual curation and
derived annotations. Give the versions used in the reported release.

\subsection{Database construction and quality control}
Explain the data model, identifier mapping, curation workflow, validation,
quality-control procedures and update process. Keep methods needed to
reproduce the work in the main manuscript.

Equation~\eqref{eq:coverage} is a notation example only. Replace it with a
justified definition and identify every term, or remove it.
\begin{equation}
 c=\frac{n_{\mathrm{included}}}{n_{\mathrm{eligible}}}.
 \label{eq:coverage}
\end{equation}

\section{Database Content}\label{sec:content}
Describe verified record types, scope and relationships. Separate measured
coverage from claims about completeness. Table~\ref{tab:content} uses
invented values to illustrate a compact inventory.

\begin{table}[t]
\centering\small
\caption{Synthetic database inventory for layout only.}
\label{tab:content}
\begin{tabularx}{\linewidth}{@{}>{\raggedright\arraybackslash}X
 S[table-format=3.0] >{\raggedright\arraybackslash}X@{}}
\toprule
Record group & {Count} & Provenance\\
\midrule
Example A & 120 & Synthetic source A\\
Example B & 240 & Synthetic source B\\
\bottomrule
\end{tabularx}
\par\smallskip\footnotesize\raggedright
These counts and source labels are invented. Replace them with the verified
content and release used in the manuscript.
\end{table}

\section{Access and Use Cases}\label{sec:access}
Give the working database URL in the body as well as the abstract.
The reserved example \url{https://example.org/database} must be replaced.
Describe querying, browsing, downloads, file formats, identifiers and any
documented programmatic access actually provided. Show a concrete use
case with interpretable results rather than only the landing page.
Figure~\ref{fig:database} is a synthetic slot for that example.

\begin{figure}[t]
\centering
\fbox{\parbox[c][33mm][c]{0.87\linewidth}{\centering
Database use-case graphic\\[5pt]Replace this synthetic slot}}
\caption{Synthetic query or use-case slot. Describe the actual query,
returned information and biological interpretation when replacing it.}
\label{fig:database}
\par\smallskip\small\raggedright
\textbf{Alt text:} Replace with a description of the actual use-case graphic.
\end{figure}

\section{Discussion and Future Maintenance}\label{sec:discussion}
Discuss verified strengths, coverage limitations, comparison with relevant
resources and the actual maintenance plan. Describe future intentions as
plans rather than completed work. Cite any supporting Supplementary
Figure~S1, Supplementary Table~S1 or Supplementary Equation~(S1) from
the separate supplement where relevant.

\section*{Data Availability}
Replace with the working resource URL, release identifier, downloadable
data and code access routes, applicable licences and any justified
restrictions. State maintenance and archiving arrangements only after
confirming them with the responsible authors.

\section*{Supplementary Data}
Describe the supplied supplementary files and replace this instruction
with the journal's required availability statement when applicable.

\section*{Author Contributions}
Replace with the contributions agreed by every author, using appropriate
CRediT roles where applicable.

\section*{Acknowledgements}
Replace with confirmed acknowledgements, or remove if not applicable.

\section*{Funding}
Replace with confirmed funder names, grant identifiers and the source of
open-access funding where applicable.

\section*{Conflict of Interest}
Replace with the authors' confirmed competing-interest disclosure.

\bibliographystyle{prism-nar}
\bibliography{references}
\end{document}
